% Options for packages loaded elsewhere
\PassOptionsToPackage{unicode}{hyperref}
\PassOptionsToPackage{hyphens}{url}
\documentclass[
  11pt,
]{article}
\usepackage{xcolor}
\usepackage[margin=1in]{geometry}
\usepackage{amsmath,amssymb}
\setcounter{secnumdepth}{-\maxdimen} % remove section numbering
\usepackage{iftex}
\ifPDFTeX
  \usepackage[T1]{fontenc}
  \usepackage[utf8]{inputenc}
  \usepackage{textcomp} % provide euro and other symbols
\else % if luatex or xetex
  \usepackage{unicode-math} % this also loads fontspec
  \defaultfontfeatures{Scale=MatchLowercase}
  \defaultfontfeatures[\rmfamily]{Ligatures=TeX,Scale=1}
\fi
\usepackage{lmodern}
\ifPDFTeX\else
  % xetex/luatex font selection
\fi
% Use upquote if available, for straight quotes in verbatim environments
\IfFileExists{upquote.sty}{\usepackage{upquote}}{}
\IfFileExists{microtype.sty}{% use microtype if available
  \usepackage[]{microtype}
  \UseMicrotypeSet[protrusion]{basicmath} % disable protrusion for tt fonts
}{}
\makeatletter
\@ifundefined{KOMAClassName}{% if non-KOMA class
  \IfFileExists{parskip.sty}{%
    \usepackage{parskip}
  }{% else
    \setlength{\parindent}{0pt}
    \setlength{\parskip}{6pt plus 2pt minus 1pt}}
}{% if KOMA class
  \KOMAoptions{parskip=half}}
\makeatother
\usepackage{longtable,booktabs,array}
\newcounter{none} % for unnumbered tables
\usepackage{calc} % for calculating minipage widths
% Correct order of tables after \paragraph or \subparagraph
\usepackage{etoolbox}
\makeatletter
\patchcmd\longtable{\par}{\if@noskipsec\mbox{}\fi\par}{}{}
\makeatother
% Allow footnotes in longtable head/foot
\IfFileExists{footnotehyper.sty}{\usepackage{footnotehyper}}{\usepackage{footnote}}
\makesavenoteenv{longtable}
\setlength{\emergencystretch}{3em} % prevent overfull lines
\providecommand{\tightlist}{%
  \setlength{\itemsep}{0pt}\setlength{\parskip}{0pt}}
\usepackage{bookmark}
\IfFileExists{xurl.sty}{\usepackage{xurl}}{} % add URL line breaks if available
\urlstyle{same}
\hypersetup{
  pdftitle={The Two-Person Rule for AI Agents: Fail-Closed Multi-Party Authorization with Offline-Verifiable Receipts},
  pdfauthor={Iman Schrock (EMILIA Protocol, Inc.; team@emiliaprotocol.ai)},
  hidelinks,
  pdfcreator={LaTeX via pandoc}}

\title{The Two-Person Rule for AI Agents: Fail-Closed Multi-Party
Authorization with Offline-Verifiable Receipts}
\author{Iman Schrock (EMILIA Protocol, Inc.; team@emiliaprotocol.ai)}
\date{20 June 2026}

\begin{document}
\maketitle

\subsection{Abstract}\label{abstract}

Autonomous agents increasingly hold credentials sufficient to perform
irreversible operations: releasing payments, changing beneficiary
records, rotating production credentials, deleting data. Existing
controls authenticate and authorize \emph{sessions and scopes}; they do
not answer the question that matters at the moment of execution ---
\emph{should this exact action happen, and which accountable humans said
yes?} We present EP-QUORUM, a multi-party authorization mechanism that
ports the two-person rule --- the control that governs nuclear release
and large-value treasury movement --- to AI-agent actions. EP-QUORUM
binds a set of pairwise-distinct, accountable human approvers, each
holding their own device-bound signing key, to one exact action, such
that the action is authorized only when a fail-closed predicate over all
of their signatures holds. The predicate enforces all-signatures-valid,
action-binding, separation of duties (distinct humans), role admission,
an M-of-N threshold, an optional total order, and a bounded approval
window. Each quorum member is an unmodified single-approver
authorization receipt over the same action hash, so the construction is
additive: a single-approver policy is the degenerate one-member case,
and the existing offline receipt verifier is reused per member. We give
the predicate, an incremental server-side admission rule that keeps a
non-conforming signer out of the trail before it is recorded, and an
adversarial conformance suite of sixteen vectors that the JavaScript,
Python, and Go reference verifiers are required to agree on. We
are deliberate about limitations: a quorum raises the cost of unilateral
action and makes every approval attributable, but it does not defeat
collusion among the required number of humans, an enrollment that lets
one human hold multiple identities, or simultaneous coercion --- and we
argue that honesty about these boundaries is a security property, not a
caveat.

\subsection{1. Introduction}\label{introduction}

The deployment of AI agents with real operational authority creates a
specific, under-addressed risk: an agent that is authenticated and
authorized \emph{in general} can nonetheless take a \emph{specific}
irreversible action that no human intended or reviewed. Identity and
access management answers ``is this actor allowed to operate in this
system?''; it does not answer ``should this particular wire, to this
beneficiary, for this amount, execute now?'' Fraud and error that occur
inside a valid session --- a prompt-injected agent, a compromised
credential, a business-email-compromise-style instruction --- are
invisible to session-level controls.

The single-approver EMILIA Protocol (EP) authorization receipt {[}1{]}
addresses this by requiring a named human to sign the exact action with
a key the orchestrating operator does not hold, producing an artifact
that is verifiable offline and consumable exactly once. But for the
highest-stakes actions, a single approver is the wrong unit of control.
The discipline that governs the most consequential decisions in finance,
government, and weapons handling is the \emph{two-person rule}: no
single individual, regardless of seniority or authentication strength,
can unilaterally cause the action; two or more distinct, accountable
people must each independently authorize.

This paper specifies EP-QUORUM, which expresses the two-person rule ---
and its generalizations to M-of-N and to ordered approval chains --- as
a cryptographic predicate over EP signoffs. Our contributions are:

\begin{enumerate}
\def\labelenumi{\arabic{enumi}.}
\tightlist
\item
  \textbf{A fail-closed multi-party authorization predicate} (Section 4)
  that composes single-approver receipts into a quorum without
  introducing any new signature type or signing ceremony, and that is
  computable offline from the receipt and policy alone.
\item
  \textbf{An incremental admission rule} (Section 5) that enforces the
  predicate as each approver signs --- rejecting a wrong-action,
  wrong-role, duplicate, out-of-order, stale, or invalid signer
  \emph{before} it enters the trail --- while requiring the executing
  system to nonetheless re-check the full predicate, because it does not
  trust the orchestrator.
\item
  \textbf{An adversarial conformance suite} (Section 6) of sixteen vectors
  carrying real WebAuthn assertions, against which the JavaScript, Python,
  and Go reference verifiers are required to agree, separating ``too strict''
  failures (denying valid quorums) from the security-critical ``too
  lenient'' failures.
\item
  \textbf{An explicit threat model and limitation statement} (Sections
  3, 7) that states what multi-party authorization does \emph{not}
  prevent.
\end{enumerate}

\subsection{2. Background and Related
Work}\label{background-and-related-work}

\textbf{Single-approver authorization receipts.} EP-QUORUM is a profile
of the EP authorization receipt {[}1{]}, which binds one approver to one
action via a user-verification-gated WebAuthn signature over a canonical
Authorization Context (action hash, policy reference, nonce, validity
window). Receipts are Merkle-anchored and verifiable fully offline, and
the protocol enforces one-time consumption and separation of duties (an
initiator must not approve its own action). Safety properties of the
single-approver state machine are maintained as machine-checked TLA+ and
Alloy models.

\textbf{Threshold and multisignature cryptography.} Threshold signatures
and multisig wallets (e.g., in distributed key generation and blockchain
custody) also require \emph{k}-of-\emph{n} parties to act. EP-QUORUM
differs in intent and verification surface: it does not aggregate keys
into one signature, and it deliberately preserves \emph{per-approver
attribution} --- the receipt records which named human signed, in which
role, at what time, over exactly which action, each independently
verifiable. The goal is accountability and offline audit, not key
aggregation; threshold cryptography optimizes the opposite (hiding the
signer set behind one key).

\textbf{Human-in-the-loop and dual control.} Dual-control and
maker-checker patterns are standard in banking software, but are
typically enforced inside a mutable application database controlled by
the very operator whose conduct an auditor would examine. EP-QUORUM
moves the evidence outside that trust boundary: each approval is a
device-held signature the operator cannot forge, and the
satisfied/not-satisfied judgment is reproducible by any third party.

\textbf{Backchannel authentication (CIBA), workload identity (WIMSE),
delegation receipts.} These authenticate the agent, transport an
approval, or bind a user's delegation; they are complementary layers and
do not produce an action-bound, offline-verifiable, multi-party-quorum
artifact. EP-QUORUM composes with them by reference.

\subsection{3. Threat Model}\label{threat-model}

We consider an agent acting in an environment with the following
adversaries.

\begin{itemize}
\tightlist
\item
  \textbf{Compromised or misaligned initiator.} The agent may be
  prompt-injected, jailbroken, or simply wrong. It can propose any
  action and any justification. We assume it cannot produce an
  approver's device-held signature.
\item
  \textbf{Compromised orchestrating operator.} The party running the
  policy registry, routing, and log may be malicious or breached. Per
  the base protocol it is \emph{not} in the signing trust path: it can
  deny service and fail to route, but cannot forge an approver's
  signature. EP-QUORUM additionally assumes the executing system
  re-verifies the quorum predicate rather than trusting the operator's
  ``approved'' assertion.
\item
  \textbf{Single compromised approver.} One approver's device,
  credential, or person may be compromised or coerced. The two-person
  rule is precisely the control that renders one such compromise
  insufficient.
\end{itemize}

We treat as \textbf{out of scope / explicitly not prevented}: collusion
among the full required set of distinct approvers; one human controlling
multiple enrolled identities (an enrollment-time control, not a runtime
one); simultaneous coercion of a complete quorum; and presentation
attacks in which an approver signs a faithful-looking but misleading
rendering (mitigated, not eliminated, by the base protocol's
rendering-fidelity controls). Section 7 states these in full.

\subsection{4. The Quorum Predicate}\label{the-quorum-predicate}

A \textbf{Quorum Policy} declares the approval \texttt{mode}
(\texttt{threshold} or \texttt{ordered}), the required count \emph{k}, a
roster of eligible \texttt{(role,\ approver)} slots, a
\texttt{distinct\_humans} flag (default true), and an approval window
\texttt{window\_sec} (default 900). A \textbf{member} is a triple
\texttt{(role,\ approver\_public\_key,\ signoff)} where the signoff is
an unmodified EP signoff --- an Authorization Context plus its WebAuthn
assertion --- over the action.

A trail of members is a \textbf{satisfied quorum} for action hash
\texttt{H} under policy \texttt{P} if and only if all of the following
hold; a verifier returns ``satisfied'' only when every check passes and
``not satisfied'' on the first failure, on a malformed policy or member,
or on any unrecognized condition:

\begin{enumerate}
\def\labelenumi{\arabic{enumi}.}
\tightlist
\item
  \textbf{Well-formed policy.} Recognized \texttt{mode}, integer
  \texttt{required\ \textgreater{}=\ 1}, non-empty roster.
\item
  \textbf{All signatures valid.} Each member's WebAuthn assertion
  verifies against its \texttt{approver\_public\_key}, with the
  assertion challenge equal to the member's context hash and user
  verification asserted. One invalid member fails the quorum.
\item
  \textbf{Action binding.} Every member's context carries
  \texttt{action\_hash\ =\ H}.
\item
  \textbf{Role admission.} Every member's \texttt{(role,\ approver)} is
  on the roster.
\item
  \textbf{Distinct humans.} Under \texttt{distinct\_humans}, approvers
  are pairwise distinct and distinct from the initiator.
\item
  \textbf{Threshold.} At least \texttt{required} admitted members exist.
\item
  \textbf{Order (ordered mode).} The \emph{i}-th member matches the
  \emph{i}-th roster slot; signature times are strictly increasing.
\item
  \textbf{Window.} Every member's \texttt{issued\_at} is within
  \texttt{window\_sec} of the first member's.
\end{enumerate}

The predicate is a pure function of \texttt{(P,\ H,\ members)}. Because
each member is an unmodified single-approver receipt, check 2 is exactly
the base verifier invoked per member; EP-QUORUM adds the set-level
checks (1, 3--8). The same function is computed by the operator before
consumption and by an independent executor or auditor offline --- they
cannot disagree on a correct input.

\textbf{Why fail-closed.} The dangerous failure in a multi-party gate is
to treat ambiguity as approval. We specify the predicate so that absence
of evidence, an unparseable member, a partial trail, or any single
failed check yields ``not satisfied,'' never ``satisfied.'' This is the
multi-party generalization of the single-approver rule that a partial or
missing signoff authorizes nothing.

\subsection{5. Incremental Admission}\label{incremental-admission}

Evaluating the full predicate only at execution time wastes the
opportunity to reject a bad signer early and produces confusing partial
trails. EP-QUORUM therefore specifies an incremental admission rule
evaluated before each new signoff is recorded: given \texttt{P},
\texttt{H}, the existing trail, and one candidate, it admits the
candidate only if the action binds, the \texttt{(role,\ approver)} is on
the roster, no admitted member shares the candidate's human (under
\texttt{distinct\_humans}), the candidate matches the next slot in
ordered mode, the time is within window (and strictly increasing in
ordered mode), and the signature verifies --- otherwise it rejects with
a specific reason (\texttt{action\_mismatch}, \texttt{ineligible\_role},
\texttt{duplicate\_human}, \texttt{out\_of\_order},
\texttt{window\_exceeded}, \texttt{non\_increasing\_time},
\texttt{invalid\_signature}).

Crucially, incremental admission is an enforcement convenience,
\textbf{not} a substitute for the predicate. The executing system
re-evaluates the full quorum gate over the assembled trail before
acting, because it does not trust the orchestrator to have admitted
honestly. This mirrors the base protocol's execution-side enforcement:
the strongest deployment places verification at the system of record,
which refuses to act on an unsatisfied quorum no matter what the
operator asserts.

\subsection{6. Conformance and
Implementation}\label{conformance-and-implementation}

We maintain an adversarial conformance suite, \texttt{EP-QUORUM-v1}, of
sixteen vectors that each carry real Class-A WebAuthn (ES256) assertions:

{\def\LTcaptype{none} % do not increment counter
\begin{longtable}[]{@{}
  >{\raggedright\arraybackslash}p{(\linewidth - 4\tabcolsep) * \real{0.2963}}
  >{\raggedright\arraybackslash}p{(\linewidth - 4\tabcolsep) * \real{0.2963}}
  >{\raggedright\arraybackslash}p{(\linewidth - 4\tabcolsep) * \real{0.4074}}@{}}
\toprule\noalign{}
\begin{minipage}[b]{\linewidth}\raggedright
Vector
\end{minipage} & \begin{minipage}[b]{\linewidth}\raggedright
Expect
\end{minipage} & \begin{minipage}[b]{\linewidth}\raggedright
Exercises
\end{minipage} \\
\midrule\noalign{}
\endhead
\bottomrule\noalign{}
\endlastfoot
\texttt{accept\_ordered\_3of3} & satisfied & Ordered PO -\textgreater{}
AO -\textgreater{} IG, distinct, increasing time \\
\texttt{accept\_threshold\_2of3} & satisfied & Any 2 distinct of 3 \\
\texttt{reject\_under\_threshold} & not satisfied & Fewer than
\texttt{required} \\
\texttt{reject\_duplicate\_human} & not satisfied & One human, two
slots \\
\texttt{reject\_out\_of\_order} & not satisfied & Out of roster order \\
\texttt{reject\_action\_mismatch} & not satisfied & Member bound to a
different action \\
\texttt{reject\_expired\_window} & not satisfied & Member outside
window \\
\texttt{reject\_one\_bad\_signature} & not satisfied & One invalid
signature \\
\texttt{reject\_wrong\_role} & not satisfied & Valid signature,
off-roster approver \\
\end{longtable}
}

The suite is run against the JavaScript, Python, and Go reference
verifiers, which are required to agree on every vector; divergence is a
conformance defect. This is a cross-language consistency check, not a
clean-room independent-implementation claim. The ``accept'' vectors guard
against an over-strict verifier (denying valid quorums); the ``reject''
vectors guard the security-critical direction (a verifier that accepts
something it must not). The mechanism is implemented in the EP reference
codebase: a pure predicate (\texttt{quorumGate}) and incremental
admission rule (\texttt{canAccept}), wired into the live
device-signature authorization path, and exercised end-to-end with
multiple virtual WebAuthn authenticators (one per approver) to confirm
that an action is blocked until the full ordered trail signs and is then
permitted.

The reuse of the single-approver verifier per member is what makes
cross-language agreement tractable: the hardest part --- WebAuthn
assertion verification and canonical context hashing --- is the base
protocol's already cross-checked code, and EP-QUORUM's additions are
set-level integer and identity checks.

\subsection{7. Limitations and Security
Considerations}\label{limitations-and-security-considerations}

We state plainly what a satisfied quorum does and does not establish.

\textbf{It proves:} \emph{k} pairwise-distinct enrolled approvers each
produced a valid, action-bound, in-window, in-order (where required)
signature with their own device-held keys, and the orchestrator forged
none of them.

\textbf{It does not prevent:} (a) \textbf{Collusion} among the required
number of distinct humans -- a \emph{k}-person rule stops \emph{k}-1 bad
actors, not \emph{k}. (b) \textbf{One human, many identities} --- if
enrollment lets a single person hold multiple approver identifiers,
distinctness is satisfied formally but not in substance; this is an
enrollment-time control (directory authority, identity proofing), out of
scope for the runtime predicate. (c) \textbf{Simultaneous coercion} of a
full quorum. (d) \textbf{Presentation attacks} --- an approver who signs
a faithful-looking but misleading rendering produces a valid signature
over the wrong understanding; the base protocol mitigates this with
rendering-fidelity and second-surface controls but cannot eliminate it
with cryptography. (e) \textbf{Diffusion of responsibility} --- adding
signers can \emph{reduce} scrutiny if each assumes another is the
careful one; quorum policies must be scoped to genuinely high-stakes,
low-frequency actions, and deployments should monitor per-role
time-to-sign and deny rates.

What EP-QUORUM converts these residual risks into is \emph{attribution}:
every approval is named, signed, role-tagged, time-stamped, and
action-bound, so an abuse is evidenced rather than deniable. We regard
the explicit statement of these boundaries as a security property ---
the most common failure in this category is a system that claims to make
fraud impossible and thereby discourages the compensating controls that
actually constrain it.

\subsection{8. Conclusion}\label{conclusion}

EP-QUORUM shows that the two-person rule, a centuries-old
human-governance control, can be expressed for AI-agent actions as a
small, fail-closed, offline-verifiable predicate over per-approver
device signatures --- additive over single-approver receipts,
cross-language conformance-tested, and honest about its limits. As
organizations move to grant agents real authority, the ability to
require that no single party --- human or machine --- can unilaterally
cause an irreversible action, while keeping every authorization
independently auditable, is a prerequisite for doing so safely.

\subsection{References}\label{references}

{[}1{]} I. Schrock, ``Authorization Receipts for High-Risk Agent Actions
(EP),'' Internet-Draft draft-schrock-ep-authorization-receipts, 2026.

{[}2{]} W3C, ``Web Authentication: An API for accessing Public Key
Credentials, Level 2.''

{[}3{]} A. Rundgren et al., ``JSON Canonicalization Scheme (JCS),'' RFC
8785, 2020.

{[}4{]} S. Bradner, ``Key words for use in RFCs to Indicate Requirement
Levels,'' BCP 14 / RFC 2119; B. Leiba, RFC 8174.

\emph{Artifacts: protocol drafts, the EP-QUORUM-v1 conformance vectors,
and the three reference implementations are available in the EMILIA
Protocol open-source repository (Apache-2.0). Site:
https://www.emiliaprotocol.ai}

\end{document}
